\subsection{2026-09-21 --- Discussion: graph structure, joint modeling, and next baseline}
\label{sec:discussion-next}

The Tokke/Vinje weak spot in Section~\ref{sec:baseline} prompted a discussion of whether the
watercourse topology (Section~\ref{sec:eda}) should be used more directly, e.g.\ via a graph
neural network over the reservoir/plant graph.

\subsubsection*{GNN considered and set aside for now.} The topology does explain the Tokke/Vinje
correlation physically (Vinje's reservoir cascades into Vinjevatn, which Tokke draws from), so a
graph-based architecture is conceptually reasonable. But: the graph is small ($\sim$22 nodes)
and \emph{identical} at train and test time, so a GNN's main advantage --- generalizing to
unseen graph structure --- doesn't really apply here. It would also be a large implementation
jump (new framework, per-case graph batching, a training loop) relative to what's been built so
far.

\subsubsection*{A premise checked and corrected.} The natural cheaper alternative --- adding each
plant's upstream-reservoir features explicitly --- turned out to already be moot: the current
modeling table (Section~\ref{sec:modeling-table}) already joins \emph{all} 17 reservoirs' data
into every row, so every unit's baseline model already had access to its neighbors' data. The
real gap is therefore not missing input features, but that the 14 units are modeled
\textbf{independently} and can't share information about each other's predictions --- a
correlated-\emph{output} problem, not a missing-\emph{input} problem.

\subsubsection*{Multi-output architectures considered.} To model 14 correlated binary outputs
without exploding the output space (treating the full ON/OFF vector as one
$2^{14}=16{,}384$-way categorical label was explicitly ruled out --- almost all combinations
never occur, and it discards the partial-credit nature of AUC scoring):
\begin{itemize}
  \item \textbf{Multi-output trees}: a single boosted model where each leaf stores a vector of
    14 probabilities instead of one scalar, with splits chosen to jointly improve all 14
    outputs (XGBoost's \texttt{multi\_strategy="multi\_output\_tree"}). Leaf count stays linear
    in tree depth, not exponential in the number of units.
  \item \textbf{Classifier chains}: feed earlier units' predicted probabilities as extra input
    features to later units' models, in some order. Simple to bolt onto the existing per-unit
    XGBoost setup.
  \item Briefly noted: a shared-trunk neural net with 14 sigmoid output heads (bigger jump from
    the current XGBoost-based approach; not pursued yet).
\end{itemize}

\subsubsection*{Decision.} Explicitly \emph{not} anchoring to the earlier plan of testing the
cheapest hypothesis first (Section~\ref{sec:eda}'s within-case price normalization) before
trying anything else --- this is an early stage of a multi-week project, and it's fine to
resequence as better questions surface. Next step agreed: build a \textbf{multi-output tree}
model as the next baseline and compare directly against the 14-independent-model baseline in
Section~\ref{sec:baseline}, using the same validation protocol (train 2015--2021, validate on
2022, score with \texttt{kaggle\_metric.py}).
